\chapter{Présentation de l'entreprise}

\section{Histoire et données clés}

Winamax est une entreprise française basée à Paris, spécialisée dans le poker et les paris sportifs en
ligne. La société trouve son origine dans un site britannique disposant d'une licence de jeu au
Royaume-Uni, racheté en 2005 par Alexandre Roos, rejoint par Christophe Schaming. Les deux
dirigeants, toujours à la tête de l'entreprise, repositionnent alors la marque autour du poker en ligne
et anticipent de plusieurs années l'ouverture du marché français.

En 2010, à la suite de l'adoption de la loi relative à l'ouverture à la concurrence et à la régulation du
secteur des jeux d'argent et de hasard en ligne, Winamax est homologuée en France pour les jeux de
cercle par l'Autorité de régulation des jeux en ligne (ARJEL), devenue l'Autorité nationale des jeux
(ANJ) depuis 2020. Le poker demeure à ce jour le seul jeu de cercle autorisé en ligne en France. En
2014, à l'approche de la Coupe du monde de football, Winamax lance son offre de paris sportifs.
L'entreprise poursuit ensuite son développement à l'international : ouverture du poker en Espagne en
2018, puis lancement plus récent des paris sportifs en Allemagne.

Aujourd'hui, Winamax revendique une communauté de plus de dix millions de joueurs répartis en
France, en Espagne et en Allemagne, et plus de deux millions de joueurs actifs. Sur le plan financier,
l'entreprise fait partie des acteurs majeurs du secteur, avec un chiffre d'affaires de plusieurs centaines
de millions d'euros. L'effectif de l'entreprise est de l'ordre de plusieurs centaines de collaborateurs,
dont une part importante travaille sur la partie technique de la plateforme.

\section{Positionnement et concurrence}

Sur le marché français du poker en ligne, Winamax occupe une position de leader avec environ la
moitié des parts de marché, portée par la plus forte liquidité du marché — plusieurs milliers de joueurs
simultanés aux heures de pointe. Plus les joueurs sont nombreux, plus l'offre de tables et de tournois
est riche, ce qui renforce l'attractivité de la plateforme. Le principal concurrent historique sur le
poker reste PokerStars, davantage implanté sur les marchés internationaux.

Sur le marché des paris sportifs, Winamax s'est hissé en tête des opérateurs français à la suite de l'Euro
de football et détient plus d'un tiers du marché, devant Betclic, Unibet et Parions Sport (FDJ). Winamax
poursuit par ailleurs une stratégie d'expansion européenne, avec une position déjà dominante sur le
poker en Espagne et une progression sur le marché allemand des paris sportifs.

Cette concurrence intense, conjuguée à la volonté de conquérir de nouveaux pays, impose à Winamax
d'innover en permanence, avec des fonctionnalités emblématiques comme l'Expresso au poker ou
MyMatch pour les paris sportifs. Ce rythme d'innovation repose directement sur la capacité de
l'équipe Infrastructure à fournir des environnements fiables et à accompagner les équipes produit
dans le déploiement de leurs services.

\section{Organisation du pôle technique}

Le pôle technique de Winamax est organisé en équipes spécialisées, chacune responsable d'un
périmètre fonctionnel précis de la plateforme :

\begin{itemize}
	\item \textbf{Client Poker} : développe et maintient le client de jeu de poker sur les différentes plateformes supportées.
	\item \textbf{Back-end Poker} : gère le moteur de jeu du poker et les transactions réalisées sur la plateforme.
	\item \textbf{Back-office Poker} : assure les interactions entre le back-end poker et l'opérateur, et met à disposition des équipes non techniques une plateforme pour, par exemple, créer des tournois.
	\item \textbf{Betting} : développe et maintient la plateforme de paris sportifs.
	\item \textbf{Betting Maths} : conçoit les algorithmes de calcul des cotes en temps réel.
	\item \textbf{Opérateur} : maintient le site web Winamax et gère le stockage des comptes joueurs.
	\item \textbf{Core} : gère les outils communs aux équipes et met en place les services exigés par les régulateurs de chaque pays.
	\item \textbf{Sécurité} : accompagne les équipes dans leurs projets et audite régulièrement la plateforme au regard des exigences des organismes de contrôle.
\end{itemize}

Cette organisation en équipes autonomes, chacune propriétaire de ses services, structure fortement
la problématique de l'alerting : lorsqu'un incident survient, il faut notifier précisément l'équipe
responsable du service réellement impacté, et non diffuser une alerte générique.

\section{L'équipe Infrastructure}

L'équipe Infrastructure est une équipe transverse et essentielle de Winamax. En lien avec le CTO, elle
décide et met en œuvre l'orientation technologique de l'entreprise, comme la migration vers AWS
décidée en 2018. Sa philosophie consiste à s'appuyer au maximum sur les services gérés fournis par
AWS et à généraliser l'Infrastructure as Code (IaC), notamment via Terraform. 

L'équipe est responsable à la fois de la mise à disposition des différents environnements, de
l'accompagnement des équipes dans le déploiement de leurs services et de la maintenance de ces
derniers. Elle prend notamment en charge plusieurs points critiques et communs à toutes les
équipes :

\begin{itemize}
	\item Provisionnement des machines nécessaires aux équipes (ECS sur EC2 ou Fargate, EC2, ECS Anywhere, machines en datacenter) ;
	\item Adressage des différentes machines et conteneurs, et filtrage des flux réseau entre les services ;
	\item Gestion de l'API Gateway (Kong) ;
	\item Déploiement des outils communs : Grafana, AutoScaler, Redmine ;
	\item Gestion de la stack d'observabilité ;
	\item Gestion des datastores pour les équipes (clusters MySQL, Redis, Kafka).
\end{itemize}

Enfin, l'équipe supervise la distribution des alertes vers les différentes équipes d'astreinte et assure
elle-même une astreinte, afin de pouvoir intervenir en cas de problème sur la plateforme. C'est
précisément dans ce périmètre que s'inscrit le sujet de mon stage : améliorer la chaîne qui relie la
détection d'un problème sur la plateforme à la notification de la bonne personne.

\section{Contexte technologique}

Historiquement, l'ensemble des technologies de Winamax était déployé sur des serveurs sur site (on-
premise). Entre 2018 et 2022, la plateforme a été majoritairement migrée vers Amazon Web Services
(AWS). Quelques serveurs demeurent toutefois dans des datacenters parisiens, hébergeant des
services liés à la régulation que la loi interdit de migrer vers le cloud. L'organisation AWS de Winamax
est découpée en de multiples comptes, un pour chaque environnement ou fonctionnalité : Production,
Prelive, Développement, PCIDSS-prod, PCIDSS-dev, Tools, Network, Keys, IT, Stacksets, Users \ldots

L'ensemble des droits IAM sur ces comptes est géré par l'équipe Sécurité, et l'accès aux
environnements se fait via des VPN OpenVPN. Côté collaboration, les équipes s'appuient sur Google
Workspace, Redmine pour les tickets, Slack pour la communication interne et GitHub pour le code
source. Slack joue un rôle central dans la problématique d'alerting, puisqu'il constitue le principal
canal de notification.

Sur le plan de l'observabilité, la plateforme combine plusieurs briques : Amazon CloudWatch, principalement pour les
métriques et les logs d'infrastructure, Amazon Managed Service for Prometheus (AMP), essentiellement pour les
métriques applicatives collectées par des collecteurs OpenTelemetry, et Grafana pour la visualisation.
C'est sur cet écosystème riche mais hétérogène que repose le système d'alerting, objet de ce rapport.

\section{Enjeux de conformité et de disponibilité}

L'activité de Winamax est soumise à une régulation stricte : l'agrément délivré par l'ANJ impose des
obligations fortes de traçabilité des opérations de jeu, de protection des joueurs et de sécurité des
données financières. C'est cette contrainte réglementaire qui explique le maintien de certains
serveurs en datacenter, hors du cloud public, ainsi que l'existence de comptes AWS dédiés et
cloisonnés à la norme PCIDSS (PCIDSS-prod, PCIDSS-dev).

Au-delà de la conformité, Winamax opère une plateforme sollicitée en continu par des millions de
joueurs, avec des pics d'activité lors des grands événements sportifs : la moindre indisponibilité se
traduit immédiatement par une dégradation de l'expérience utilisateur, une perte de revenus et un
risque de réputation. C'est dans ce contexte d'exigence réglementaire et opérationnelle que s'inscrit
le rôle de l'équipe Infrastructure, et la nécessité d'un système d'alerting fiable, capable de détecter
rapidement tout incident affectant la disponibilité ou la conformité de la plateforme.
